\chapter{从基因到蛋白质}\label{ch:g11:gene-expression}

1961 年，一位生物化学家向细菌的无\omterm{def:g10:cells-common-unit:cell}{细胞}提取液中加入
人工合成的 \omterm{def:g11:gene-expression:transcription}{RNA}，这种 \omterm{def:g11:gene-expression:transcription}{RNA} 仅由重复的字母 U 组成。该提取液
合成出由单一\omterm{def:g10:chemistry-of-life:families}{氨基酸}苯丙氨酸重复连接而成的\omterm{def:g11:gene-expression:protein}{蛋白质}。
这一实验破译了\omterm{def:g11:gene-expression:code}{遗传密码}的第一个\omterm{def:g11:gene-expression:code}{密码子}：UUU 编码苯丙氨酸。
五年之内，全部 64 个\omterm{def:g11:gene-expression:code}{密码子}均已破译，
而且这些\omterm{def:g11:gene-expression:code}{密码子}在细菌、小麦和人体内
竟然完全相同。本章沿着从\omterm{def:g10:universal-dna:gene}{基因}的\omterm{def:g10:universal-dna:nucleotide}{核苷酸}序列
到\omterm{def:g11:gene-expression:protein}{蛋白质}\omterm{def:g10:chemistry-of-life:families}{氨基酸}序列的途径展开——这正是
\cref{ch:g11:mutations}中每一种性状所遵循的路径。

\section{基因制造蛋白质}

\begin{proposition}[一个基因，一种蛋白质]\label{prop:g11:gene-expression:onegene}
\omterm{def:g10:universal-dna:gene}{基因}被表达，是指\omterm{def:g10:cells-common-unit:cell}{细胞}利用其序列
合成它所编码的\omterm{def:g11:gene-expression:protein}{蛋白质}。大多数性状依赖于\omterm{def:g11:gene-expression:protein}{蛋白质}——酶、
结构纤维、载体、受体——而\omterm{def:g11:mutations:mutation}{突变}通过改变\omterm{def:g10:universal-dna:gene}{基因}所产生的
\omterm{def:g11:gene-expression:protein}{蛋白质}来改变性状。
\end{proposition}
\begin{proof}[证据]
比德尔和塔特姆（1941 年）用射线照射面包霉的孢子，
筛选出在基本培养基上不能生长、只有添加某种特定物质
（如\omterm{def:g10:chemistry-of-life:families}{氨基酸}精氨酸）才能生长的\omterm{def:g11:mutations:mutation}{突变}体。
每个\omterm{def:g11:mutations:mutation}{突变}体都缺失了合成该物质的反应链中的一种酶；
每种缺陷都由单个\omterm{def:g10:universal-dna:gene}{基因}遗传；处于同一反应链不同步骤上
受阻的\omterm{def:g11:mutations:mutation}{突变}体则携带位于不同\omterm{def:g10:universal-dna:gene}{基因}中的\omterm{def:g11:mutations:mutation}{突变}。
一个\omterm{def:g10:universal-dna:gene}{基因}，一种酶——后来的研究将其推广为
一个\omterm{def:g10:universal-dna:gene}{基因}，一条多肽链。
\end{proof}

\begin{definition}[蛋白质，氨基酸序列]\label{def:g11:gene-expression:protein}
\emph{蛋白质}\index{蛋白质}是由\omterm{def:g10:chemistry-of-life:families}{氨基酸}连接而成的链，长度从几十个到
数千个不等，所用\omterm{def:g10:chemistry-of-life:families}{氨基酸}取自二十种。
\emph{序列}——即\omterm{def:g10:chemistry-of-life:families}{氨基酸}的排列顺序——由\omterm{def:g10:universal-dna:gene}{基因}所规定；
一旦合成，肽链便折叠成确定的三维结构，
该结构由序列决定，并进一步决定功能。
改变一个\omterm{def:g10:chemistry-of-life:families}{氨基酸}，折叠方式可能改变，
功能也可能随之改变。
\end{definition}

\section{转录：将基因复制为 RNA}

\begin{definition}[信使 RNA 与转录]\label{def:g11:gene-expression:transcription}
\emph{RNA}\index{RNA}（核糖\omterm{def:g10:chemistry-of-life:families}{核酸}）是单链的\omterm{def:g10:universal-dna:nucleotide}{核苷酸}链，
与 \omterm{def:g10:universal-dna:information}{DNA} 的\omterm{def:g10:universal-dna:nucleotide}{核苷酸}相似，区别在于其糖为核糖，
碱基尿嘧啶（U）代替胸腺嘧啶（T）。\emph{转录}\index{转录}
是将\omterm{def:g10:universal-dna:gene}{基因}序列复制为一个
\emph{信使 RNA}\index{信使 RNA}（mRNA）分子的过程：酶
\emph{RNA 聚合酶}沿\omterm{def:g10:universal-dna:gene}{基因}打开\omterm{prop:g10:universal-dna:helix}{双螺旋}，读取一条
链（即\emph{模板链}），并以其为模板组装互补的
RNA，U 与 A 配对，A 与 T 配对，G 与 C 配对，C 与 G 配对。mRNA
的序列与\omterm{def:g10:universal-dna:gene}{基因}的另一条链相同，只是 U 代替 T；它离开
\omterm{def:g10:cells-common-unit:organelle}{细胞核}进入\omterm{def:g10:cells-common-unit:cell}{细胞质}。
\end{definition}

\begin{omfigure}
\begin{tikzpicture}[font=\small, >=Stealth, line cap=round]
  % DNA strands with bubble
  \draw[very thick, omDef] (0,0.5) -- (2.5,0.5) .. controls (3.2,0.5) and (3.4,1.3) .. (4.3,1.3) -- (6.7,1.3) .. controls (7.6,1.3) and (7.8,0.5) .. (8.5,0.5) -- (11,0.5);
  \draw[very thick, omDef] (0,-0.5) -- (2.5,-0.5) .. controls (3.2,-0.5) and (3.4,-1.3) .. (4.3,-1.3) -- (6.7,-1.3) .. controls (7.6,-1.3) and (7.8,-0.5) .. (8.5,-0.5) -- (11,-0.5);
  \foreach \x in {0.3,0.7,...,2.3} \draw[gray] (\x,0.5) -- (\x,-0.5);
  \foreach \x in {8.7,9.1,...,10.9} \draw[gray] (\x,0.5) -- (\x,-0.5);
  % template strand letters (bottom strand in bubble)
  \foreach \x/\l in {4.5/T, 4.9/A, 5.3/C, 5.7/G, 6.1/G, 6.5/A} \node[font=\scriptsize, omDef, anchor=north] at (\x,-1.32) {\l};
  % mRNA
  \draw[very thick, omThm] (4.4,-0.9) -- (6.6,-0.9) .. controls (7.4,-0.9) and (7.2,0.1) .. (8.0,0.6) .. controls (8.6,1.0) and (9.4,1.4) .. (10.6,1.7);
  \foreach \x/\l in {4.5/A, 4.9/U, 5.3/G, 5.7/C, 6.1/C, 6.5/U} \node[font=\scriptsize, omThm, anchor=south] at (\x,-0.85) {\l};
  \foreach \x in {4.5,4.9,...,6.5} \draw[gray, densely dotted] (\x,-0.95) -- (\x,-1.28);
  % polymerase
  \draw[thick, fill=omProp!30, rounded corners=8pt] (6.95,-1.6) rectangle (8.65,0.2);
  \node[align=center] at (7.8,-0.7) {RNA\\ 聚合酶};
  \draw[->, thick] (8.7,-1.9) -- (9.8,-1.9) node[right] {沿基因移动};
  \node[anchor=east, omDef, align=right] at (-0.15,0) {基因的 \\ DNA};
  \node[anchor=west, omThm, align=left] at (10.7,1.7) {正在释放的 \\ mRNA};
  \node[anchor=north, omDef] at (5.5,-1.75) {模板链};
\end{tikzpicture}

{\small \omterm{def:g11:gene-expression:transcription}{转录}。\omterm{def:g11:gene-expression:transcription}{RNA} 聚合酶打开\omterm{prop:g10:universal-dna:helix}{双螺旋}，读取
模板链，并合成与之互补的信使 \omterm{def:g11:gene-expression:transcription}{RNA}（以 U 代替 T）。
酶通过后，\omterm{prop:g10:universal-dna:helix}{双螺旋}重新闭合，
新合成的 \omterm{def:g11:gene-expression:transcription}{RNA} 随之脱落。}
\end{omfigure}

\begin{example}[一条转录本]\label{ex:g11:gene-expression:transcript}
模板链 \texttt{3'-TAC GGA CTT ATC-5'} \omterm{def:g11:gene-expression:transcription}{转录}出 mRNA
\texttt{5'-AUG CCU GAA UAG-3'}。\omterm{def:g10:universal-dna:gene}{基因}的另一条链为
\texttt{5'-ATG CCT GAA TAG-3'}：mRNA 的序列与该链相同，
只是将 T 换为 U，因此\omterm{def:g10:universal-dna:gene}{基因}序列通常按
这条链书写，即\emph{编码链}。
\end{example}

\section{遗传密码}

\begin{definition}[密码子与遗传密码]\label{def:g11:gene-expression:code}
mRNA 按连续的三个\omterm{def:g10:universal-dna:nucleotide}{核苷酸}为一组读取，这些组称为
\emph{密码子}\index{密码子}，读取从确定的起点开始。
\emph{遗传密码}\index{遗传密码}是 64 种可能的密码子与
20 种\omterm{def:g10:chemistry-of-life:families}{氨基酸}之间的对应关系：61 个密码子各编码一种\omterm{def:g10:chemistry-of-life:families}{氨基酸}，
三个密码子（UAA、UAG、UGA）是\emph{终止}信号，用于结束
肽链；AUG 编码甲硫氨酸，同时也作为\emph{起始}密码子。
该遗传密码是\emph{简并}的——多数\omterm{def:g10:chemistry-of-life:families}{氨基酸}对应
多个密码子——并且是\emph{通用}的：同一套密码表
适用于所有已知生物，仅有个别罕见的微小变异。
\end{definition}

\begin{omfigure}
\begin{tabular}{c|cccc|c}
\multicolumn{1}{c}{} & \multicolumn{4}{c}{\small 第二个字母} & \multicolumn{1}{c}{} \\
{\small 第一个字母} & U & C & A & G & {\small 第三个字母} \\ \hline
U & UUU Phe & UCU Ser & UAU Tyr & UGU Cys & U \\
 & UUC Phe & UCC Ser & UAC Tyr & UGC Cys & C \\
 & UUA Leu & UCA Ser & UAA \textbf{终止} & UGA \textbf{终止} & A \\
 & UUG Leu & UCG Ser & UAG \textbf{终止} & UGG Trp & G \\ \hline
C & CUU Leu & CCU Pro & CAU His & CGU Arg & U \\
 & CUC Leu & CCC Pro & CAC His & CGC Arg & C \\
 & CUA Leu & CCA Pro & CAA Gln & CGA Arg & A \\
 & CUG Leu & CCG Pro & CAG Gln & CGG Arg & G \\ \hline
A & AUU Ile & ACU Thr & AAU Asn & AGU Ser & U \\
 & AUC Ile & ACC Thr & AAC Asn & AGC Ser & C \\
 & AUA Ile & ACA Thr & AAA Lys & AGA Arg & A \\
 & AUG \textbf{Met} & ACG Thr & AAG Lys & AGG Arg & G \\ \hline
G & GUU Val & GCU Ala & GAU Asp & GGU Gly & U \\
 & GUC Val & GCC Ala & GAC Asp & GGC Gly & C \\
 & GUA Val & GCA Ala & GAA Glu & GGA Gly & A \\
 & GUG Val & GCG Ala & GAG Glu & GGG Gly & G \\
\end{tabular}

{\small \omterm{def:g11:gene-expression:code}{遗传密码}，在 mRNA 上读取。每个\omterm{def:g11:gene-expression:code}{密码子}由其
第一个字母（行）、第二个字母（列）和第三个字母（块内的
行）确定。AUG 既编码甲硫氨酸，又是起始信号；三个
\omterm{def:g11:gene-expression:code}{密码子}为终止\omterm{def:g11:gene-expression:code}{密码子}。}
\end{omfigure}

\begin{proposition}[遗传密码的破译]\label{prop:g11:gene-expression:readcode}
每个\omterm{def:g11:gene-expression:code}{密码子}的含义都是通过实验确定的。
\end{proposition}
\begin{proof}[证据]
尼伦伯格和马太（1961 年）在无\omterm{def:g10:cells-common-unit:cell}{细胞}提取物中加入由单一
重复\omterm{def:g10:universal-dna:nucleotide}{核苷酸}组成的人工合成 \omterm{def:g11:gene-expression:transcription}{RNA}，该提取物含有\omterm{def:g10:cells-common-unit:organelle}{核糖体}、\omterm{def:g10:chemistry-of-life:families}{氨基酸}
以及\omterm{def:g11:gene-expression:protein}{蛋白质}合成所需的酶。poly-U 生成苯丙氨酸链，poly-C 生成
脯氨酸链，poly-A 生成赖氨酸链。
随后，用重复二联体和三联体的人工合成 \omterm{def:g11:gene-expression:transcription}{RNA}，以及
单个三联体与\omterm{def:g10:cells-common-unit:organelle}{核糖体}结合的实验，为其余\omterm{def:g11:gene-expression:code}{密码子}作了归属，至
1966 年完成。三字母长度此前已由算术推出
（两个字母只能给出 16 种组合，不足以对应 20 种\omterm{def:g10:chemistry-of-life:families}{氨基酸}；
三个字母可给出 64 种），并由\omterm{def:g11:mutations:mutation}{突变}证实：在\omterm{def:g10:universal-dna:gene}{基因}中插入一个或两个
\omterm{def:g10:universal-dna:nucleotide}{核苷酸}会破坏其\omterm{def:g11:gene-expression:protein}{蛋白质}，插入三个则
可恢复出一个几乎正常的\omterm{def:g11:gene-expression:protein}{蛋白质}。
\end{proof}

\begin{method}[序列的翻译]\label{met:g11:gene-expression:translate}
\begin{enumerate}
  \item 写出编码链（或\omterm{def:g11:gene-expression:transcription}{转录}模板链）：
        将 T 替换为 U，即得 mRNA。
  \item 找到第一个 AUG：它即起始\omterm{def:g11:gene-expression:code}{密码子}，阅读框由此
        确定。
  \item 从 AUG 起将 mRNA 切分为连续三联体，并逐一
        查表。
  \item 遇到第一个终止\omterm{def:g11:gene-expression:code}{密码子}即停止，依次列出的\omterm{def:g10:chemistry-of-life:families}{氨基酸}
        即为\omterm{def:g11:gene-expression:protein}{蛋白质}的序列。
  \item 对\omterm{def:g11:mutations:mutation}{突变}体，重复上述步骤并比较：\omterm{def:g11:gene-expression:protein}{蛋白质}不变（沉默\omterm{def:g11:mutations:mutation}{突变}），一个
        \omterm{def:g10:chemistry-of-life:families}{氨基酸}改变（错义\omterm{def:g11:mutations:mutation}{突变}），提前终止（无义\omterm{def:g11:mutations:mutation}{突变}），
        \omterm{def:g11:mutations:mutation}{突变}之后的全部序列改变（移码\omterm{def:g11:mutations:mutation}{突变}）。
\end{enumerate}
\end{method}

\begin{example}[四种突变，四种结果]\label{ex:g11:gene-expression:outcomes}
编码链 \texttt{ATG CCT GAA TTC TAG}：mRNA
\texttt{AUG CCU GAA UUC UAG}，\omterm{def:g11:gene-expression:protein}{蛋白质}为 Met--Pro--Glu--Phe。
\begin{itemize}
  \item \texttt{ATG CC\textbf{C} GAA TTC TAG}：CCC 仍为脯氨酸——
        \emph{沉默\omterm{def:g11:mutations:mutation}{突变}}。
  \item \texttt{ATG CCT G\textbf{C}A TTC TAG}：GCA 为丙氨酸——
        Met--Pro--Ala--Phe，\emph{错义\omterm{def:g11:mutations:mutation}{突变}}。
  \item \texttt{ATG CCT \textbf{T}AA TTC TAG}：UAA 为终止\omterm{def:g11:gene-expression:code}{密码子}——
        Met--Pro，\emph{无义\omterm{def:g11:mutations:mutation}{突变}}，得到截短的\omterm{def:g11:gene-expression:protein}{蛋白质}。
  \item \texttt{ATG CC\textbf{A}T GAA TTC TAG}：插入的 A 使阅读框
        移位：\texttt{AUG CCA UGA}\dots{} Met--Pro--终止——
        \emph{移码\omterm{def:g11:mutations:mutation}{突变}}。
\end{itemize}
\end{example}

\section{翻译：读取信息}

\begin{proposition}[翻译]\label{prop:g11:gene-expression:translation}
\emph{翻译}\index{翻译}发生在
\omterm{def:g10:cells-common-unit:cell}{细胞质}中的\emph{\omterm{def:g10:cells-common-unit:organelle}{核糖体}}上。\omterm{def:g10:cells-common-unit:organelle}{核糖体}在 mRNA 的
起始\omterm{def:g11:gene-expression:code}{密码子}处与之结合，并沿其逐个\omterm{def:g11:gene-expression:code}{密码子}移动。每个\omterm{def:g11:gene-expression:code}{密码子}都由一个
\emph{转运 \omterm{def:g11:gene-expression:transcription}{RNA}}（tRNA）与之配对。tRNA 是一种小分子 \omterm{def:g11:gene-expression:transcription}{RNA}，其一端
携带与\omterm{def:g11:gene-expression:code}{密码子}互补的三\omterm{def:g10:universal-dna:nucleotide}{核苷酸}\emph{反\omterm{def:g11:gene-expression:code}{密码子}}，另一端携带
相应的\omterm{def:g10:chemistry-of-life:families}{氨基酸}。\omterm{def:g10:cells-common-unit:organelle}{核糖体}将每次带来的\omterm{def:g10:chemistry-of-life:families}{氨基酸}
接到正在延伸的肽链上，并释放空载的 tRNA；
遇到终止\omterm{def:g11:gene-expression:code}{密码子}时，释放已完成的肽链，肽链随即折叠。数个
\omterm{def:g10:cells-common-unit:organelle}{核糖体}依次阅读同一条 mRNA，一条 mRNA 可产生数百
个该\omterm{def:g11:gene-expression:protein}{蛋白质}分子，之后该 mRNA 被降解。
\end{proposition}
\admitted

\begin{omfigure}
\begin{tikzpicture}[font=\small, >=Stealth, line cap=round]
  % mRNA
  \draw[very thick, omThm] (0,0) -- (11,0);
  \foreach \x/\l in {0.4/A,0.8/U,1.2/G, 1.7/C,2.1/C,2.5/U, 3.0/G,3.4/A,3.8/A, 4.3/U,4.7/U,5.1/C, 5.6/G,6.0/G,6.4/U, 6.9/U,7.3/A,7.7/G}
    \node[font=\scriptsize, omThm, anchor=north] at (\x,-0.05) {\l};
  \node[anchor=east, omThm] at (-0.15,0) {信使 RNA};
  % ribosome (two lobes) over codons 3 and 4
  \draw[thick, fill=omProp!25, rounded corners=10pt] (2.7,0.15) rectangle (5.5,1.4);
  \draw[thick, fill=omProp!35, rounded corners=12pt] (2.5,1.4) .. controls (2.5,3.0) and (5.7,3.0) .. (5.7,1.4) -- cycle;
  \node at (4.1,2.1) {核糖体};
  % tRNAs
  \draw[very thick, omMeth!80!black] (3.4,0.55) -- (3.4,1.35);
  \node[font=\scriptsize, omMeth!80!black, anchor=south] at (3.4,0.15) {CUU};
  \draw[very thick, omMeth!80!black] (4.7,0.55) -- (4.7,1.35);
  \node[font=\scriptsize, omMeth!80!black, anchor=south] at (4.7,0.15) {AAG};
  % amino acids chain
  \foreach \i/\x in {0/3.4, 1/4.7} \fill[omDef!70] (\x,1.5) circle (4pt);
  \foreach \x in {2.6, 1.9, 1.2} \fill[omDef!70] (\x,1.9) circle (4pt);
  \draw[thick, omDef!70] (1.2,1.9) -- (2.6,1.9) -- (3.4,1.5) -- (4.7,1.5);
  \node[anchor=east, omDef!80!black, align=right] at (0.9,1.9) {新生 \\ 蛋白质};
  % incoming tRNA
  \draw[very thick, omMeth!80!black] (7.0,1.3) -- (7.0,2.1);
  \fill[omDef!70] (7.0,2.25) circle (4pt);
  \node[font=\scriptsize, omMeth!80!black, anchor=west] at (7.15,1.35) {CCA};
  \node[anchor=west, align=left, font=\scriptsize] at (7.3,1.7) {下一个 tRNA 及其 \\ 所携氨基酸，即将到达};
  \draw[->, thick] (7.0,1.1) -- (6.1,0.75);
  \draw[->, thick] (5.7,-0.9) -- (7.2,-0.9) node[right] {核糖体移动};
  \node[anchor=north, font=\scriptsize] at (4.05,-0.45) {密码子};
  \draw (2.85,-0.35) -- (2.85,-0.45) -- (3.95,-0.45) -- (3.95,-0.35);
\end{tikzpicture}

{\small \omterm{prop:g11:gene-expression:translation}{翻译}。\omterm{def:g10:cells-common-unit:organelle}{核糖体}同时容纳两个\omterm{def:g11:gene-expression:code}{密码子}，与各\omterm{def:g11:gene-expression:code}{密码子}配对的
反\omterm{def:g11:gene-expression:code}{密码子}所在的 tRNA 携带其\omterm{def:g10:chemistry-of-life:families}{氨基酸}，肽链转移至新到的\omterm{def:g10:chemistry-of-life:families}{氨基酸}上，
\omterm{def:g10:cells-common-unit:organelle}{核糖体}前移一个
\omterm{def:g11:gene-expression:code}{密码子}。遇到终止\omterm{def:g11:gene-expression:code}{密码子}时，肽链被释放。}
\end{omfigure}

\begin{example}[速度与产量]\label{ex:g11:gene-expression:speed}
一个\omterm{def:g10:cells-common-unit:organelle}{核糖体}每秒可加入 5 至 20 个\omterm{def:g10:chemistry-of-life:families}{氨基酸}：一条含 300 个\omterm{def:g10:chemistry-of-life:families}{氨基
酸}的\omterm{def:g11:gene-expression:protein}{蛋白质}用时不到一分钟。\omterm{def:g10:cells-common-unit:organelle}{核糖体}相继排列在同一条 mRNA 上，
彼此相距约 80 个\omterm{def:g10:universal-dna:nucleotide}{核苷酸}，因此一条含 900 个\omterm{def:g10:universal-dna:nucleotide}{核苷酸}的 mRNA 可同时承载
十二个\omterm{def:g10:cells-common-unit:organelle}{核糖体}；在其几小时的寿命内，一条 mRNA 可产生
一千个\omterm{def:g11:gene-expression:protein}{蛋白质}分子。一个 \emph{E.~coli} \omterm{def:g10:cells-common-unit:cell}{细胞}约
每秒制造 $\num{20000}$ 个\omterm{def:g11:gene-expression:protein}{蛋白质}分子；一个红\omterm{def:g10:cells-common-unit:cell}{细胞}的
前体一生中约产生 $\num{5e8}$ 个血红蛋白分子。
\end{example}

\section{一个基因，多种信息}

\begin{proposition}[真核生物中 mRNA 的成熟]\label{prop:g11:gene-expression:splicing}
在\omterm{def:g10:cells-common-unit:prokaryote}{真核细胞}中，\omterm{def:g10:universal-dna:gene}{基因}序列被
\emph{内含子}\index{内含子}打断，\omterm{prop:g11:gene-expression:splicing}{内含子}是被\omterm{def:g11:gene-expression:transcription}{转录}、但随后
在 \omterm{def:g11:gene-expression:transcription}{RNA} 离开\omterm{def:g10:cells-common-unit:organelle}{细胞核}前被切除的片段；保留并首尾相接的片段，
即\emph{外显子}，构成成熟的 mRNA。许多\omterm{def:g10:universal-dna:gene}{基因}
可以有多种剪切方式，保留不同组合的外显子
（\emph{可变剪接}），因此一个\omterm{def:g10:universal-dna:gene}{基因}可产生多种
不同的 mRNA 以及若干相关的\omterm{def:g11:gene-expression:protein}{蛋白质}：$\num{20000}$ 个人类\omterm{def:g10:universal-dna:gene}{基因}
可编码远超 $\num{100000}$ 种\omterm{def:g11:gene-expression:protein}{蛋白质}。
\end{proposition}
\admitted

\begin{omfigure}
\begin{tikzpicture}[font=\small, >=Stealth]
  \tikzset{ex/.style={draw, fill=omDef!40, minimum height=0.45cm, minimum width=1.0cm, inner sep=1pt},
           intr/.style={draw, fill=gray!20, minimum height=0.45cm, minimum width=0.7cm, inner sep=1pt}}
  \node[anchor=east] at (-0.2,0) {基因（DNA）};
  \node[ex] at (0.5,0) {1}; \node[intr] at (1.35,0) {}; \node[ex] at (2.2,0) {2}; \node[intr] at (3.05,0) {}; \node[ex] at (3.9,0) {3}; \node[intr] at (4.75,0) {}; \node[ex] at (5.6,0) {4};
  \node[anchor=west, font=\scriptsize] at (6.2,0) {外显子（蓝色），内含子（灰色）};
  \draw[->, thick] (3.0,-0.35) -- (3.0,-0.95) node[midway, right] {转录};
  \node[anchor=east] at (-0.2,-1.3) {前体 mRNA};
  \node[ex, fill=omThm!30] at (0.5,-1.3) {1}; \node[intr] at (1.35,-1.3) {}; \node[ex, fill=omThm!30] at (2.2,-1.3) {2}; \node[intr] at (3.05,-1.3) {}; \node[ex, fill=omThm!30] at (3.9,-1.3) {3}; \node[intr] at (4.75,-1.3) {}; \node[ex, fill=omThm!30] at (5.6,-1.3) {4};
  \draw[->, thick] (2.0,-1.65) -- (1.0,-2.35) node[midway, left] {剪接};
  \draw[->, thick] (4.0,-1.65) -- (6.0,-2.35) node[midway, right] {可变剪接};
  \node[anchor=east] at (-0.6,-2.7) {成熟 mRNA A};
  \node[ex, fill=omThm!30] at (0.0,-2.7) {1}; \node[ex, fill=omThm!30] at (1.0,-2.7) {2}; \node[ex, fill=omThm!30] at (2.0,-2.7) {3}; \node[ex, fill=omThm!30] at (3.0,-2.7) {4};
  \node[anchor=west] at (3.6,-2.7) {蛋白质 A};
  \node[anchor=east] at (5.4,-2.7) {};
  \node[ex, fill=omThm!30] at (6.0,-2.7) {1}; \node[ex, fill=omThm!30] at (7.0,-2.7) {2}; \node[ex, fill=omThm!30] at (8.0,-2.7) {4};
  \node[anchor=west] at (8.6,-2.7) {蛋白质 B};
\end{tikzpicture}

{\small 真核\omterm{def:g10:universal-dna:gene}{基因}先被整体\omterm{def:g11:gene-expression:transcription}{转录}，随后\omterm{prop:g11:gene-expression:splicing}{内含子}被
切除。保留全部外显子或舍去其中一个，同一\omterm{def:g10:universal-dna:gene}{基因}便得到两种 mRNA
并由此得到两种\omterm{def:g11:gene-expression:protein}{蛋白质}。}
\end{omfigure}

\begin{remark}[信息的流向]\label{rem:g11:gene-expression:flow}
\omterm{def:g10:universal-dna:information}{DNA} \omterm{def:g11:gene-expression:transcription}{转录}为 \omterm{def:g11:gene-expression:transcription}{RNA}，\omterm{def:g11:gene-expression:transcription}{RNA} \omterm{prop:g11:gene-expression:translation}{翻译}为\omterm{def:g11:gene-expression:protein}{蛋白质}，\omterm{def:g11:gene-expression:protein}{蛋白质}
决定性状：信息只沿一个方向流动。改变了的
\omterm{def:g11:gene-expression:protein}{蛋白质}从不改写\omterm{def:g10:universal-dna:gene}{基因}，这正是一生的训练
或晒伤无法遗传的原因；只有 \omterm{def:g10:universal-dna:information}{DNA} 的\omterm{def:g11:mutations:mutation}{突变}
（见\cref{ch:g11:mutations}）才会改变下一代得到的\omterm{def:g10:universal-dna:information}{遗传信息}。
每个生物体的每个\omterm{def:g10:cells-common-unit:cell}{细胞}中都进行着同样的三个步骤，
使用的是同一套\omterm{def:g11:gene-expression:code}{遗传密码}：\cref{ch:g10:universal-dna}所述的普遍性，
如今已细化至最细微之处。
\end{remark}

\section{练习}

\begin{exercise}[$\star$]\label{exo:g11:gene-expression:1}
列出 \omterm{def:g10:universal-dna:information}{DNA} 与 \omterm{def:g11:gene-expression:transcription}{RNA} 之间的三点不同。
\end{exercise}

\begin{exercise}[$\star$]\label{exo:g11:gene-expression:2}
写出以模板链 \texttt{3'-TAC CGA AAA ATT-5'} \omterm{def:g11:gene-expression:transcription}{转录}得到的
mRNA。
\end{exercise}

\begin{exercise}[$\star$]\label{exo:g11:gene-expression:3}
借助\omterm{def:g11:gene-expression:code}{遗传密码}表，\omterm{prop:g11:gene-expression:translation}{翻译} mRNA \texttt{AUG GGC AAA UGU UAA}。
\end{exercise}

\begin{exercise}[$\star$]\label{exo:g11:gene-expression:4}
在\omterm{prop:g11:gene-expression:translation}{翻译}过程中，\omterm{def:g10:cells-common-unit:organelle}{核糖体}与转运 \omterm{def:g11:gene-expression:transcription}{RNA}
各起什么作用？
\end{exercise}

\begin{exercise}[$\star$]\label{exo:g11:gene-expression:5}
为什么\omterm{def:g11:gene-expression:code}{遗传密码}每个\omterm{def:g10:chemistry-of-life:families}{氨基酸}至少需要三个\omterm{def:g10:universal-dna:nucleotide}{核苷酸}？
\end{exercise}

\begin{exercise}[$\star\star$]\label{exo:g11:gene-expression:6}
一条\omterm{def:g11:gene-expression:protein}{蛋白质}含 412 个\omterm{def:g10:chemistry-of-life:families}{氨基酸}。求其 mRNA 编码部分（含终止\omterm{def:g11:gene-expression:code}{密码子}）
的最短长度，以及对应\omterm{def:g10:universal-dna:gene}{基因}
的最短长度（以碱基对计）。
\end{exercise}

\begin{exercise}[$\star\star$]\label{exo:g11:gene-expression:7}
编码链 \texttt{ATG AAA GGC TGG TAA} \omterm{def:g11:mutations:mutation}{突变}为
\texttt{ATG AAA GGC TGA TAA}。写出两种\omterm{def:g11:gene-expression:protein}{蛋白质}，并判断该
\omterm{def:g11:mutations:mutation}{突变}的类型。
\end{exercise}

\begin{exercise}[$\star\star$]\label{exo:g11:gene-expression:8}
起始序列相同，分别\omterm{def:g11:mutations:mutation}{突变}为 \texttt{ATG AAG GGC TGG TAA}
和 \texttt{ATG AAA GGG CTG GTA A}。判断每种\omterm{def:g11:mutations:mutation}{突变}的类型，并写出对应的\omterm{def:g11:gene-expression:protein}{蛋白质}。
\end{exercise}

\begin{exercise}[$\star\star$]\label{exo:g11:gene-expression:9}
解释为何\omterm{def:g11:gene-expression:code}{密码子}第三位的替换往往
不改变所编码的\omterm{def:g10:chemistry-of-life:families}{氨基酸}，并结合表格说明。
\end{exercise}

\begin{exercise}[$\star\star$]\label{exo:g11:gene-expression:10}
多聚 UC \omterm{def:g11:gene-expression:transcription}{RNA}（\texttt{UCUCUCUC}\dots）\omterm{prop:g11:gene-expression:translation}{翻译}出的\omterm{def:g11:gene-expression:protein}{蛋白质}
由丝氨酸与亮氨酸交替组成。证明这与三联体密码相符，
并据此确定两个\omterm{def:g11:gene-expression:code}{密码子}。
\end{exercise}

\begin{exercise}[$\star\star$]\label{exo:g11:gene-expression:11}
一个人类\omterm{def:g10:universal-dna:gene}{基因}长 \num{30000} 个碱基对，产生一条含 500
个\omterm{def:g10:chemistry-of-life:families}{氨基酸}的\omterm{def:g11:gene-expression:protein}{蛋白质}。解释这一差异的原因。
\end{exercise}

\begin{exercise}[$\star\star\star$]\label{exo:g11:gene-expression:12}
某药物能抑制细菌的\omterm{def:g10:cells-common-unit:organelle}{核糖体}而不抑制人类的\omterm{def:g10:cells-common-unit:organelle}{核糖体}。
解释其为何可作抗生素，并说明这一事实表明两类生物的\omterm{def:g10:cells-common-unit:organelle}{核糖体}
存在怎样的差异。
\end{exercise}

\begin{exercise}[$\star\star\star$]\label{exo:g11:gene-expression:13}
解释\cref{ch:g10:universal-dna}中的水母\omterm{def:g10:universal-dna:gene}{基因}如何
被小鼠读取（用本章的术语），以及
若小鼠使用另一套密码，会发生什么。
\end{exercise}

\begin{exercise}[$\star\star\star$]\label{exo:g11:gene-expression:14}
一次\omterm{def:g11:mutations:mutation}{突变}改变了反\omterm{def:g11:gene-expression:code}{密码子}为 UAG 的 tRNA，使其
携带一个\omterm{def:g10:chemistry-of-life:families}{氨基酸}，而不是使\omterm{prop:g11:gene-expression:translation}{翻译}终止。预测其对\omterm{def:g10:cells-common-unit:cell}{细胞}的
\omterm{def:g11:gene-expression:protein}{蛋白质}的总体影响，以及对携带提前出现的
UAG 的\omterm{def:g11:mutations:mutation}{突变}\omterm{def:g10:universal-dna:gene}{基因}的影响。
\end{exercise}

\begin{exercise}[$\star\star\star$]\label{exo:g11:gene-expression:15}
讨论为何\omterm{def:g10:universal-dna:gene}{基因}起始端附近的移码\omterm{def:g11:mutations:mutation}{突变}几乎总比
末端附近的移码\omterm{def:g11:mutations:mutation}{突变}更有害，并讨论错义\omterm{def:g11:mutations:mutation}{突变}为何
可以从无害到致命不一。
\end{exercise}

\section{问题：一个基因的四种读法}

\begin{problem}[{周末问题——一个短基因的转录、
翻译、三次突变与计时：从 DNA 到氨基酸，
以及一个蛋白质所需的六十秒}]
\label{pb:g11:gene-expression:1}
一个短\omterm{def:g10:universal-dna:gene}{基因}的编码链序列为：
\begin{center}
\texttt{ATG GCT TGG AAA CCC GAA TAC GGT CAT TTC TGA}
\end{center}

\medskip
\noindent\textbf{第一部分——读取\omterm{def:g10:universal-dna:gene}{基因}。}
\begin{enumerate}
  \item 写出模板链。
  \item 写出由该\omterm{def:g10:universal-dna:gene}{基因}\omterm{def:g11:gene-expression:transcription}{转录}得到的 mRNA。
  \item 逐个\omterm{def:g11:gene-expression:code}{密码子}进行\omterm{prop:g11:gene-expression:translation}{翻译}，
        给出\omterm{def:g11:gene-expression:protein}{蛋白质}的\omterm{def:g10:chemistry-of-life:families}{氨基酸}序列。
  \item 该\omterm{def:g11:gene-expression:protein}{蛋白质}含有多少个\omterm{def:g10:chemistry-of-life:families}{氨基酸}，
        编码这些\omterm{def:g10:chemistry-of-life:families}{氨基酸}需用多少个 mRNA \omterm{def:g10:universal-dna:nucleotide}{核苷酸}，
        包括终止\omterm{def:g11:gene-expression:code}{密码子}？
  \item 哪一个\omterm{def:g11:gene-expression:code}{密码子}起始了\omterm{prop:g11:gene-expression:translation}{翻译}，
        又是什么确定了所有其他\omterm{def:g11:gene-expression:code}{密码子}的阅读框？
\end{enumerate}

\medskip
\noindent\textbf{第二部分——三种\omterm{def:g11:mutations:mutation}{突变}体。}
\begin{enumerate}[resume]
  \item \omterm{def:g11:mutations:mutation}{突变}体 1：\texttt{ATG GCT TGG AAA CCG GAA TAC GGT CAT TTC TGA}。
        给出\omterm{def:g11:gene-expression:protein}{蛋白质}，并对\omterm{def:g11:mutations:mutation}{突变}进行分类。
  \item \omterm{def:g11:mutations:mutation}{突变}体 2：\texttt{ATG GCT TGG AAA CCC GAA TAC GGT CAT TTC TGA}，
        其第 9 个\omterm{def:g11:gene-expression:code}{密码子} \texttt{CAT} 被替换为 \texttt{CGT}。
        给出\omterm{def:g11:gene-expression:protein}{蛋白质}并进行分类。
  \item \omterm{def:g11:mutations:mutation}{突变}体 3：\texttt{ATG GCT TGA AAA CCC GAA TAC GGT CAT TTC TGA}。
        给出\omterm{def:g11:gene-expression:protein}{蛋白质}并进行分类。哪一次单个替换
        造成了这一\omterm{def:g11:mutations:mutation}{突变}？
  \item \omterm{def:g11:mutations:mutation}{突变}体 4：在第 6 个\omterm{def:g10:universal-dna:nucleotide}{核苷酸}之后插入一个 T：
        \texttt{ATG GCT TTG GAA ACC CGA ATA CGG TCA TTT CTG A}。给出
        \omterm{def:g11:gene-expression:protein}{蛋白质}并进行分类。
  \item 将四个\omterm{def:g11:mutations:mutation}{突变}体按对\omterm{def:g11:gene-expression:protein}{蛋白质}的破坏程度
        由小到大排序，并说明理由。
\end{enumerate}

\medskip
\noindent\textbf{第三部分——\omterm{def:g11:gene-expression:code}{遗传密码}为何如此。}
\begin{enumerate}[resume]
  \item 有多少个\omterm{def:g11:gene-expression:code}{密码子}编码亮氨酸？丝氨酸？色氨酸？哪一种
        \omterm{def:g10:chemistry-of-life:families}{氨基酸}的\omterm{def:g11:gene-expression:code}{密码子}最易因随机替换而改变，
        哪一种最不易？
  \item 计算所有发生在第三位上的替换中，
        亮氨酸\omterm{def:g11:gene-expression:code}{密码子} \texttt{CUx} 第三位上属沉默的比例。
  \item 说明为何两个碱基的密码行不通，
        以及为何不需要四个碱基的密码。
  \item 多聚 U 产生多聚苯丙氨酸，多聚 A 产生多聚赖氨酸。
        这两个实验确立了密码表中的哪些条目？
  \item \omterm{def:g11:gene-expression:code}{遗传密码}在细菌和人类中相同。请说明
        其对\omterm{prop:g10:universal-dna:universal}{转基因技术}的意义，
        以及对亲缘关系的意义。
\end{enumerate}

\medskip
\noindent\textbf{第四部分——为“\omterm{def:g11:gene-expression:protein}{蛋白质}工厂”计时。}
一个\omterm{def:g10:cells-common-unit:organelle}{核糖体}每秒加入 10 个\omterm{def:g10:chemistry-of-life:families}{氨基酸}，\omterm{def:g10:cells-common-unit:organelle}{核糖体}在 mRNA 上依次跟随，
彼此相隔 80 个\omterm{def:g10:universal-dna:nucleotide}{核苷酸}；每条 mRNA 的寿命为 2 小时。
\begin{enumerate}[resume]
  \item 一个\omterm{def:g10:cells-common-unit:organelle}{核糖体}\omterm{prop:g11:gene-expression:translation}{翻译}该\omterm{def:g10:universal-dna:gene}{基因}编码的\omterm{def:g11:gene-expression:protein}{蛋白质}需要多长时间，
        那么一个含 600 个\omterm{def:g10:chemistry-of-life:families}{氨基酸}的\omterm{def:g11:gene-expression:protein}{蛋白质}呢？
  \item 一条 1800 个\omterm{def:g10:universal-dna:nucleotide}{核苷酸}的 mRNA 上，
        同时可有多少个\omterm{def:g10:cells-common-unit:organelle}{核糖体}在读取？
  \item 若每 8 秒起始一个新的\omterm{def:g10:cells-common-unit:organelle}{核糖体}，
        一条 mRNA 在其寿命内能产生多少个该\omterm{def:g11:gene-expression:protein}{蛋白质}的拷贝？
  \item 某细菌需要在 10 分钟内获得 \num{2000} 个
        该酶的拷贝。同时\omterm{def:g11:gene-expression:transcription}{转录}该\omterm{def:g10:universal-dna:gene}{基因}，需要多少条
        mRNA？
  \item 说明结果：该\omterm{def:g10:universal-dna:gene}{基因}编码的\omterm{def:g11:gene-expression:protein}{蛋白质}、
        四个\omterm{def:g11:mutations:mutation}{突变}中使其失活的那一个，以及一个
        \omterm{def:g10:cells-common-unit:organelle}{核糖体}合成它所需的时间。
\end{enumerate}
\end{problem}
